% Einheit 4 von 5 – Sachaufgaben (bank/lineare-gleichungssysteme/e4.jsonl, zusammenbau v0.1)
%% TODO Zweigzeile Teil 1: was hier gelernt wird (Satz in Schülersprache aus dem Einheitstitel) – Einheit 4
\einheitenkopf[e4]{Einheit 4 von 5 · Sachaufgaben}
\zweigzeile{ab Kl. 9 (am Gymnasium ab Kl. 8) · P10}
\setcounter{aufgabe}{22}

%% TODO \verfahren{…}: Verfahrensüberschrift in Sachsprache fehlt in der Bank (2.3 b)
%% TODO Titel: Ich-kann-Satz fehlt in der Bank (Platzhalter: Kettenname) – Nr. 23
%% TODO Anweisung: ein Satz über der ersten Teilaufgabe fehlt in der Bank – Nr. 23
\begin{aufgabe}{Sachaufgaben}
%% TODO \rechenplatz{n}: Schrittzahl des Grundfalls fehlt in der Bank (nur bei fester Schreibform, 2.3 b)
\begin{teile}
\teil Bäckerei: $3$ Brötchen und $2$ Brezeln kosten $3{,}10\,€$; $1$ Brötchen und $4$ Brezeln kosten $3{,}70\,€$. Was ist unbekannt? Benenne die gesuchten Größen mit $x$ und $y$ und schreib heraus, welche Zahlen rechts vom Gleichheitszeichen und welche vor $x$ und $y$ stehen.
\teil Ein Eis und ein Kakao kosten zusammen $4{,}60\,€$. Drei Eis und vier Kakao kosten $15{,}80\,€$. $x$: Preis eines Eises, $y$: Preis eines Kakaos (in €). Gleichungen? Stelle I und II auf. I: \leerfeld \quad II: \leerfeld
\teil Freibad: Am Morgen wurden $50$ Karten verkauft, Kinderkarten zu $2{,}80\,€$ und Erwachsenenkarten zu $4{,}50\,€$, zusammen für $177{,}40\,€$. $x$: Anzahl der Kinderkarten, $y$: Anzahl der Erwachsenenkarten. Gleichungen? Stelle I und II auf. I: \leerfeld \quad II: \leerfeld
\teil Jugendherberge: $23$ Zimmer, Zweibett- und Vierbettzimmer, mit zusammen $68$ Betten. $x$: Anzahl der Zweibettzimmer, $y$: Anzahl der Vierbettzimmer. Gleichungen? Stelle I und II auf. I: \leerfeld \quad II: \leerfeld
\teil Die Summe zweier Zahlen ist $31$, ihre Differenz ist $7$. $x$: die größere Zahl, $y$: die kleinere Zahl. Gleichungen? Stelle I und II auf. I: \leerfeld \quad II: \leerfeld
\teil Schreibwaren: Ein Block und ein Füller kosten zusammen $8\,€$, vier Blöcke und zwei Füller $26\,€$. Dazu gehört I: $x + y = 8$, II: $4x + 2y = 26$. Wofür stehen $x$ und $y$?
\teil Obstkorb: $x$ ist die Anzahl der Äpfel, $y$ die Anzahl der Birnen im Korb. Was sagt $x + y = 25$? Formuliere einen Satz.
\end{teile}
\begin{gleichungsraster}[2]
\gl{\text{Imbiss: Zwei Burger und eine Portion Pommes kosten $14{,}50\,€$, ein Burger und zwei Portionen Pommes $12{,}50\,€$. Was kostet ein Burger, was eine Portion Pommes? Stelle ein System auf, löse es und antworte.}} &
\end{gleichungsraster}
\begin{teile}
\teil Café: Ein Croissant und ein Muffin kosten zusammen $4{,}10\,€$. Fünf Croissants und zwei Muffins kosten $13{,}60\,€$. $x$ ist der Preis eines Croissants, $y$ der Preis eines Muffins in Euro. Stelle ein Gleichungssystem auf. \hfill (P10 2024 OS) \\ I: \leerfeld \quad II: \leerfeld
\end{teile}
\end{aufgabe}

%% TODO \verfahren{…}: Verfahrensüberschrift in Sachsprache fehlt in der Bank (2.3 b)
%% TODO Titel: Ich-kann-Satz fehlt in der Bank (Platzhalter: Kettenname) – Nr. 24
%% TODO Anweisung: ein Satz über der ersten Teilaufgabe fehlt in der Bank – Nr. 24
\begin{aufgabe}{Sachaufgaben, Sek II}
%% TODO \rechenplatz{n}: Schrittzahl des Grundfalls fehlt in der Bank (nur bei fester Schreibform, 2.3 b)
\begin{teile}
\teil Was ist unbekannt? „Ein Rechteck hat den Umfang $30$ cm; die eine Seite $x$ ist $3$ cm länger als die andere.“ Kreuze an, wofür $x$ steht.\\ \kreuz{eine Länge}\\ \kreuz{eine Anzahl}\\ \kreuz{ein Preis}\\ \kreuz{ein Anteil}
\teil Bäckerei: $x$ ist die Zahl der Roggenbrote, $y$ die Zahl der Weizenbrote. I: $x + y = 60$, II: $x = 2y$. Was sagen I und II? Formuliere je einen Satz.
\teil Tierfutter: Drei Sorten enthalten $15\,\%$, $25\,\%$ und $45\,\%$ Eiweiß. Eine Mischung aus den Anteilen $u$, $v$ und $w$ der drei Sorten soll $30\,\%$ Eiweiß haben: I: $u + v + w = 1$, II: $15u + 25v + 45w = 30$. Was bedeuten I und II?
\teil Aus drei Düngern mit $4\,\%$, $8\,\%$ und $16\,\%$ Stickstoff wird ein Dünger mit $9\,\%$ Stickstoff gemischt. Dazu gehört I: $4u + 8v + 16w = 9$, II: $u + v + w = 1$. Interpretiere das Gleichungssystem im Sachzusammenhang. \hfill (Abitur 2024 LK)
\end{teile}
\end{aufgabe}

%% TODO Titel: Ich-kann-Satz fehlt in der Bank (Platzhalter: Kettenname) – Nr. 25
%% TODO Anweisung: ein Satz über der ersten Teilaufgabe fehlt in der Bank – Nr. 25
\begin{aufgabe}{Mischungen und Alter (Vorrat)}
\begin{gleichungsraster}[2]
\gl{\text{Studentenfutter: Nüsse zu $14\,€$ je kg und Rosinen zu $5\,€$ je kg sollen $30$ kg Mischung zu $8\,€$ je kg ergeben. Wie viel kg von jeder Sorte? Stelle ein System auf und löse es.}} & \\
\end{gleichungsraster}
\end{aufgabe}

%% TODO Titel: Ich-kann-Satz fehlt in der Bank (Platzhalter: Kettenname) – Nr. 26
%% TODO Anweisung: ein Satz über der ersten Teilaufgabe fehlt in der Bank – Nr. 26
\begin{aufgabe}{Probe im Text (Rechnung mit den Zahlen der Aufgabe)}
\begin{teile}
\teil Probe im Text: „Eine Tasse und ein Teller kosten zusammen $9\,€$; drei Tassen und zwei Teller kosten $23\,€$.“ Lisa findet: Tasse $5\,€$, Teller $4\,€$. Prüfe mit den Zahlen aus dem Text. \janein
\end{teile}
\end{aufgabe}

%% TODO Titel: Ich-kann-Satz fehlt in der Bank (Platzhalter: Kettenname) – Nr. 27
%% TODO Anweisung: ein Satz über der ersten Teilaufgabe fehlt in der Bank – Nr. 27
\begin{aufgabe}{Bestandteile einer Gleichung deuten}
\begin{teile}
\teil Handyrechnung: $x$ ist die Zahl der Minuten, $y$ die Zahl der SMS. Was bedeutet in $0{,}09x + 0{,}19y = 12{,}60$ die Zahl $0{,}09$, was die Zahl $12{,}60$?
\end{teile}
\end{aufgabe}

%% TODO Titel: Ich-kann-Satz fehlt in der Bank (Platzhalter: Kettenname) – Nr. 28
%% TODO Anweisung: ein Satz über der ersten Teilaufgabe fehlt in der Bank – Nr. 28
\begin{aufgabe}{Sachaufgaben – Fehler finden, Begründen}
\begin{teile}
\teil „Ein Kakao und ein Muffin kosten zusammen $4{,}20\,€$; zwei Kakao und fünf Muffins kosten $14{,}10\,€$.“ Ole stellt auf: I: $x + y = 14{,}10$, II: $2x + 5y = 4{,}20$. Finde den Fehler.
\teil „Ein Hotel hat $17$ Zimmer, Zweibett- und Dreibettzimmer, mit zusammen $43$ Betten.“ $x$ ist die Anzahl der Zweibettzimmer, $y$ die Anzahl der Dreibettzimmer. Kim rechnet $x = 8$, $y = 9$ und antwortet: „Es gibt $8$ Dreibettzimmer und $9$ Zweibettzimmer.“ Finde den Fehler.
\teil Müsli: Drei Sorten mit $5\,\%$, $15\,\%$ und $30\,\%$ Nussanteil werden zu einem Müsli mit $25\,\%$ Nussanteil gemischt: I: $u + v + w = 1$, II: $5u + 15v + 30w = 25$. Nils schreibt: „$u$, $v$ und $w$ sind die Nussanteile der drei Sorten in Prozent.“ Finde den Fehler.
\teil Warum reicht die Angabe „Ein Brot und ein Kuchen kosten zusammen $8\,€$“ nicht, um beide Preise zu bestimmen? Begründe.
\teil Eine Mischung aus drei Sorten hat die Anteile $u$, $v$ und $w$. Warum gilt $u + v + w = 1$, und warum ist die Gleichung für den Gehalt einer Zutat eine Bilanz? Begründe.
\end{teile}
\end{aufgabe}
